\documentclass[openany,fontset=fandol]{ctexbook}
\usepackage{amsmath,amssymb}
\usepackage[unicode,hidelinks]{hyperref}
\setlength{\emergencystretch}{2em}
\interfootnotelinepenalty=10000
\DeclareMathOperator{\Spec}{Spec}
\DeclareMathOperator{\Aut}{Aut}
\DeclareMathOperator{\Div}{Div}
\newcommand{\OO}{\mathcal O}
\newcommand{\PP}{\mathbf P}
\raggedbottom
\begin{document}
\title{代数几何基础第五卷预备笔记\\参数丛单态性与射影标架}
\author{据法文扫描和英译编辑稿互校续译}
\date{2026年11月22日}
\maketitle
\chapter*{范围与接续}
第83批收录第三篇命题4.3的陈述、反例、域底情形的证明及所用的射影标架引理。
法文物理140--143页（手写147--150），英译物理106--108页。
下一处是法文物理143页、英译物理108页的任意底一般情形，本批不将它计为已译或已证明。
正式出版的 EGA IV.1 另行接续；本批尚未并入1158页合订本。

\chapter*{第三篇（续）}
\section*{4\quad 除子线性系与可逆模（续）}

\noindent\textbf{命题4.3。}
设 \(D\) 是 \(X/S\) 上由 \(P^\vee\) 参数化的除子线性系，其中 \(g:X\to S\) 平坦且局部有限表示。\footnote{法140在“有限表示”上方补写“局部”；(a)的边注另要求 \(g\) 有限表示。英106把两层条件混在一起。本命题的总假设取局部有限表示，(a)再加强，(b)只另要求满射。}

\noindent\textbf{a)}
再设 \(g\) 有限表示。对每个 \(s\in S\)，令 \(k\) 为 \(k(s)\) 的一个代数闭包，且假设几何纤维 \(X_k\) 有一个关联素循环 \(T\)，使自然映射
\[
 k\longrightarrow H^0(T,\OO_T)
\]
是同构。如果 \(g\) 固有且满射，这一条件自动满足。于是在此情形下，态射
\[
 D:P^\vee\longrightarrow\Div_{X/S}
\]
是单态射。\footnote{这里 \(T\) 是几何纤维的某个关联点的闭包，随后赋予约化结构；不要求该点是唯一的关联点。固有满射时纤维非空，选取一个不可约分支的泛点；其约化闭包在代数闭域上固有且整，故全局函数为该域，见 \href{https://stacks.math.columbia.edu/tag/04L0}{Stacks，04L0}。法141写 \(\Div_{X/S}\)，英106误写 \(\Div_{S/S}\)。}

\noindent\textbf{b)}
考察映射 \(u\mapsto D\circ u\)，它把 \(\Aut_S(P^\vee)\) 送入 \(X/S\) 上由 \(P^\vee\) 参数化的除子族的集合。如果 \(g\) 满射，这个映射为单射；特别地，\(D=D\circ u\) 蕴含 \(u=\mathrm{id}_{P^\vee}\)。更一般地，自同构函子到以 \(P^\vee\) 参数化的除子线性系函子的态射是单态射。\footnote{最后一句据法141手写补入；它指每次换底后的映射也为单射，而不是只比较 \(S\) 的点。用 \(u\) 重参数化仍保持线性系。}

在 a) 的假设下，b) 是 a) 的直接推论；但 b) 只需要较弱的假设。应当注意，若去掉 a) 的额外条件，结论可以不成立。例如取 \(S=\Spec k\)，取 \(\PP^1_k\) 的一个非空开子概形 \(X\)，不含两个不同的 \(k\)-点 \(a,b\)。令 \(D\) 为 \(X\hookrightarrow\PP^1_k\) 所定义的除子族。则 \(a,b\) 在 \(X\) 上定义同一除子，即零除子，但 \(a\ne b\)。\footnote{法141明确给出定义除子族的典范浸入；英106省略这一句。此处添明 \(X\ne\varnothing\) 以用非平凡反例。例如取 \(X=\PP^1_k\setminus\{a,b\}\)，这两个点除子分别限制到 \(X\) 均为空，仍是有效 Cartier 除子。}

\medskip\noindent\textbf{域底情形。}
先设 \(S\) 是域 \(k\) 的谱。利用下降，可假定 \(k\) 代数闭。选取 a) 中的 \(T\)，赋予诱导的约化结构。把所得线性系的除子限制到 \(T\)，只须证明相应复合
\[
 P^\vee\longrightarrow\Div_{T/k}
\]
为单态射。限制所得的族仍是除子线性系，因此归结到 \(X=T\)，从而归结到 \(k\xrightarrow{\sim} H^0(X,\OO_X)\) 的情形。\footnote{法141和英107都把这一步略写成 \(\Div_{X/k}\to\Div_{T/k}\)；任意除子未必能这样限制，例如在 \(\mathbb A^2_k\) 上把坐标直线作为除子限制到它自身，其方程恒零。因此正文只写这个线性系的限制。\(T\) 的关联泛点处，正则截面不落入那个关联素理想；对几何点同样成立。故在约化整 \(T\) 上得到非零截面，再用相对 Cartier 判据，见补证。}

对任意 \(k\)-概形 \(S'\)，此时有 \(\OO_{S'}\xrightarrow{\sim}g_{S',*}\OO_{X_{S'}}\)。
设 \(L\) 和 \(u:E\to g_*L\) 如4.1。如果 \(E_{S'}\) 的两条处处不为零的截面 \(\phi,\psi\) 所给出的 \(L_{S'}\) 截面 \(u(\phi),u(\psi)\) 有相同的除子，那么两截面相差 \(\OO_{X_{S'}}\) 的一个可逆整体截面。这个单位来自 \(S'\)；于是 \(\phi,\psi\) 也只相差底上的一个单位，从而给出同一个 \(P^\vee(S')\) 的点。每个 \(P^\vee(S')\) 的点在 \(S'\) 局部都由 \(E_{S'}\) 的一条处处不为零的截面给出（见第一章），这就证明了域底情形的 a)。\footnote{英107在 \(u\) 的靶处保留问号，按4.1补成 \(g_*L\)。这里用 \(S'\) 区别检验概形和原底 \(\Spec k\)。最后一步还用评估映射对参数截面的单射性；补证给出有限仿射开检验，覆盖非约化的 \(S'\)。}

\medskip\noindent\textbf{引理。}
设 \(k\) 代数闭，\(X\) 是非空、在 \(k\) 上局部有限型的概形，\(D\) 是 \(X\) 上由 \(P^\vee=\PP(E^\vee)\) 参数化的除子线性系。考虑相应的态射
\[
 f:X-Z\longrightarrow P=\PP(E),
\]
其中 \(Z\) 是基点集。若 \(r=\dim_k E>0\)，则可取 \(r+1\) 个点 \(x_i\in X(k)-Z(k)\)，\(1\leq i\leq r+1\)，使 \(f(x_i)\) 构成 \(P\) 的一个射影标架：任取 \(r\) 个指标，其对应的像不全在一个超平面内。\footnote{法142手写补齐 \(x_i\) 和指标范围，英107保留。秩一时 \(P\) 只有一个点、没有射影超平面，标架条件无实质内容，可重用同一点；不可额外要求 \(X\) 在秩一时一定有两个不同的 \(k\)-点。}

可以把 \(X\) 换成 \(X-Z\)，因而假定 \(Z=\varnothing\)。由4.1，\(f(X)\) 不在任何超平面内，因此可先选出 \(r\) 个点 \(x_1,\ldots,x_r\)，其像射影独立，也就是由 \(E\) 上的 \(r\) 个线性无关的线性形式给出。
剩下的是选 \(x_{r+1}\in X(k)\)，使其像避开从这 \(r\) 个像中各取 \(r-1\) 个所张成的 \(r\) 个超平面 \(H_i\)。否则由雅各布森性质可知
\[
 f(X)\subseteq\bigcup_i H_i .
\]
于是，若 \(X_0\) 是 \(X\) 的任一不可约分支，\(f(X_0)\) 必落入某个 \(H_i\)，与4.1矛盾。\footnote{法142划去“先把 \(X\) 换成一个整分支”的长段。有效论证是：非零参数截面在各分支的关联泛点均不消失，所以逐分支非退化；仅有 \(f(X)\) 整体非退化，不足以得到射影标架。闭点稠密用 \href{https://stacks.math.columbia.edu/tag/0477}{Stacks，0477}。}

引理既得，为证明域底情形的 b)，也可假定 \(Z=\varnothing\)。由4.1，只须证明知道 \(u\) 在 \(P\) 上的逆转置（逆步）自同构与 \(f:X\to P\) 的复合，就能确定 \(P^\vee\) 的自同构 \(u\)；并且对任意基变换 \(S'\to\Spec k\)，同样能确定 \(P^\vee_{S'}\) 的相对自同构。由刚才的引理和第1小节对 \(\PP(E^\vee)\) 的自同构的描述，这一点随即成立：射影丛的自同构由它在一组射影标架截面上的作用确定，所谓标架截面，是指它们在各纤维上构成射影标架。\footnote{法142明确是逆转置（逆步）自同构与 \(f\) 的复合；采用 \(\PP(E)\) 表示一维商的约定。推导时用的是一组在底上定义的截面，不是只知道自同构在每个剩余域上的约化。见下方矩阵核查。}

\section*{译校补证（非原文）}
\paragraph{限制与底上的单位。}
\(T\) 是代数闭域上的约化整有限型概形，因此几何整。对任一几何参数，原族的截面在 \(T\) 的关联泛点按原纤维正则性不消失；具体地，在含泛点的仿射邻域 \(\Spec A\) 中，关联素理想 \(\mathfrak p\) 给出嵌入 \(A/\mathfrak p\hookrightarrow A\)。扩域后嵌入仍成立，且 \((A/\mathfrak p)\otimes_k K\) 为整环，故扩域后的该泛点仍为关联点。在 \(T\) 的几何纤维上，线丛的非零截面正则。于是该族在 \(T\) 上仍有相对有效除子；限制参数方向的 \(\OO(1)\) 不变。比较原族中相等的除子，可通过方程限制到 \(T\) 后继续比较。
对 \(k\) 上有限表示的 \(T\)，取有限仿射覆盖和其交的有限仿射覆盖；计算全局函数所用的等化子与任何平坦系数扩张交换。对 \(S'\) 的仿射开 \(\Spec R\)，\(R\) 是平坦 \(k\)-模，故 \(H^0(T_R,\OO)=R\)，单位也随之对应。

\paragraph{从相同除子回到参数。}
在 \(T\) 上，\(k\)-线性评估 \(E\to H^0(T,L|_T)\) 单射，因为非零向量的截面正则。把截面限制到 \(T\) 的若干仿射开。由于 \(E\) 有限维，可选有限多个开，使各限制映射的核之交仍为零。对任一 \(S'\) 的仿射开，将合映射与其坐标环作 \(k\)-张量，仍为单射；它通过 \(u\) 在换底后对参数的评估。若 \(u(\psi)=a\,u(\phi)\) 且 \(a\) 是底上的单位，则 \(\psi=a\phi\)，包括 \(S'\) 非约化的情形。相同 Cartier 除子的局部单位比值唯一，因而粘成整体单位。对非代数闭 \(k\) 用忠实平坦的扩域检测两个参数态射是否相等。

\paragraph{标架在任意换底上检测自同构。}
把域上的标架变到坐标点及点 \([1:\cdots:1]\)，再把这些点拉回为 \(S'\) 上的截面。
一个固定所有标架截面的射影自同构，局部可用可逆矩阵表示（用 \(\OO(1)\) 的拉回与底线丛的局部平凡化）。固定各坐标截面迫使矩阵的全部非对角元为零。再固定最后一个截面，所有对角元相等，故该自同构为恒等。这里相等是环上的等式，不是只查看每个点的剩余域；例如在 \(k[\epsilon]/(\epsilon^2)\) 上，矩阵 \(I+\epsilon E_{12}\) 约化为 \(I\)，对应的射影自同构却不恒等。命题4.3的任意底情形在底本下一段另证。
\end{document}
